\documentclass[../../../include/open-logic-section]{subfiles}

\begin{document}

\olfileid[pt]{sth}{z}{pairs}
\olsection{Pares}

O próximo axioma a considerar é o seguinte:

\begin{axiom}[Pares]
Quaisquer que sejam os conjuntos $a, b$, o conjunto $\{a, b\}$ existe.
\[
	\forall a \forall b \exists P \forall x (x \in P \liff (x = a \lor x = b))
\]
\end{axiom}

Eis como justificar este axioma à luz da concepção iterativa. Suponhamos que $a$ está disponível no estágio $S$, e $b$ no estágio $T$. Seja $M$ o que quer que seja dos estágios $S$ e $T$ que ocorre mais tarde. Então, uma vez que $a$ e $b$ estão ambos disponíveis no estágio $M$, o conjunto $\{a,b\}$ é uma coleção possível disponível em qualquer estágio depois de $M$ (o que quer que seja o posterior).

Mas espera!{} Porquê assumir que \emph{existem} estágios depois de $M$? Se não existir nenhum, a nossa justificação falha. Por isso, para justificar Pares, teremos de acrescentar outro princípio à história que contamos na \olref[sth][z][story]{sec}, nomeadamente:
\begin{enumerate}
	\item[] \stagessucc. Não há último estágio.
\end{enumerate}
Este princípio está justificado? Nada na história de Shoenfield afirmava
\emph{explicitamente} que não há último estágio. Ainda assim, mesmo que seja
(estritamente falando) um acréscimo à nossa história, encaixa bem na
ideia básica de que os conjuntos se formam em estágios. Aceitá-lo-emos simplesmente no que se segue. E, assim, aceitaremos também o Axioma dos Pares.

Armados com este novo axioma, podemos provar a existência de muitos mais conjuntos. Por exemplo:

\begin{prop}\ollabel{prop:pairsconsequences}
Quaisquer que sejam os conjuntos $a$ e $b$, existem os seguintes conjuntos:
	\begin{enumerate}
		\item\ollabel{singleton} $\{a\}$
		\item\ollabel{binunion} $a \cup b$
		\item\ollabel{tuples} $\tuple{a, b}$
	\end{enumerate}
\end{prop}

\begin{proof}
\olref{singleton}. Por Pares, $\{a, a\}$ existe, o que é $\{a\}$ por
Extensionalidade.

\olref{binunion}. Por Pares, $\{a, b\}$ existe. Ora, $a \cup b = \bigcup
\{a, b\}$ existe por União.

\olref{tuples}. Por \olref{singleton}, $\{a\}$ existe. Por Pares, $\{a,
b\}$ existe. Ora, $\{\{a\}, \{a, b\}\} = \tuple{a, b}$ existe, de novo por Pares.
\end{proof}

\begin{prob}
Mostre que, quaisquer que sejam os conjuntos $a, b, c$, o conjunto $\{a, b, c\}$ existe.
\end{prob}

\begin{prob}
Mostre que, quaisquer que sejam os conjuntos $a_1, \ldots, a_n$, o conjunto $\{a_1, \ldots,
a_n\}$ existe.
\end{prob}

%\begin{proof} By two applications of Pairs, $\{\{a_1, a_2\}, \{a_1,
%   a_3\}\}$ exists. By Union and Extensionality, $\bigcup \{\{a_1,
%   a_2\}, \{a_1, a_3\}\} = \{a_1, a_2, a_3\}$ exists. Repeat this
%   trick as often as necessary. \end{proof}

\end{document}
